% 全統記述レベル 実戦模試 数学 標準編 第1回（見本版メタデータ）
\newcommand{\ExamTitle}{全統記述レベル 実戦模試 数学}
\newcommand{\ExamSubtitle}{標準編 第1回}
\newcommand{\ExamShort}{実戦模試 標準編 第1回}
\newcommand{\ExamScope}{数学I・II・A・B・C}
\newcommand{\ExamTime}{100}
\newcommand{\ExamFullMark}{100}
\newcommand{\ExamDaimonNum}{4}
\newcommand{\ExamYear}{2027}
\newcommand{\ExamSampleBadge}{見本（大問1のみ収録・製品版は全4題）}


\input{store-preamble}
\renewcommand{\BookletPart}{解答・解説}
\hypersetup{
  pdftitle={\ExamTitle\ \ExamSubtitle{}｜解答・解説},
  pdfauthor={大学別数学研究室／森 祐太},
  pdfsubject={\ExamYear 年度入試対策 \ExamScope}
}

\begin{document}

\storecover{解\ \ 答\ \ ・\ \ 解\ \ 説}{%
  \item[1.] 本冊子には，試験場で解答用紙にそのまま書き写せる\textbf{論述形式の解答}を小問の順に掲載している。
  \item[2.] 二重下線を引いた「\textbf{答：}」が各小問の最終結論である。そこに辿り着くまでの場合分け・増減表・根拠が自分の答案に書けているかを確かめること。
  \item[3.] 灰色の帯で示した「\textbf{別解}」や「\textbf{確かめ}」は，復習および検算の手引きとして活用すること。
  \item[4.] 自己採点の配点（加点・減点）は，別冊「採点基準」を参照すること。
}

\section*{解答・解説}

\hdA{大問 1 ── 3次関数の極値・実数解の個数・区間における最小値}
\hdTag{分野：微分法（数学II）\quad 配点：25点（(1) 8点／(2) 7点／(3) 10点）}

\Q{1}{8}
$f(x)=x^{3}-3ax^{2}+4$ を $x$ で微分すると
\[
  f'(x)=3x^{2}-6ax=3x(x-2a)
\]
である。$a>0$ であるから $0<2a$ であり，$f'(x)=0$ となるのは $x=0,\ 2a$ のときである。
よって，$f(x)$ の増減表は次のようになる。

\begin{center}\small
\renewcommand{\arraystretch}{1.35}
\begin{tabular}{|c||c|c|c|c|c|}
\hline
$x$ & $\cdots$ & $0$ & $\cdots$ & $2a$ & $\cdots$ \\
\hline
$f'(x)$ & $+$ & $0$ & $-$ & $0$ & $+$ \\
\hline
$f(x)$ & $\nearrow$ & 極大 & $\searrow$ & 極小 & $\nearrow$ \\
\hline
\end{tabular}
\end{center}

ここで
\[
  f(0)=4,\qquad
  f(2a)=(2a)^{3}-3a(2a)^{2}+4=8a^{3}-12a^{3}+4=4-4a^{3}
\]
であるから，$f(x)$ は $x=0$ で極大値 $4$，$x=2a$ で極小値 $4-4a^{3}$ をとる。

\Ans{\text{$x=0$ で極大値 $4$，\quad $x=2a$ で極小値 $4-4a^{3}$}}

\Q{2}{7}
3次関数 $y=f(x)$ のグラフは，(1) の増減表より $x\leqq0$ で単調に増加し，
$0\leqq x\leqq2a$ で単調に減少し，$x\geqq2a$ で単調に増加する。
したがって，方程式 $f(x)=0$ が異なる3つの実数解をもつための必要十分条件は，
曲線 $y=f(x)$ と $x$ 軸が異なる3つの共有点をもつこと，すなわち
\[
  \text{（極大値）}>0\quad\text{かつ}\quad\text{（極小値）}<0
\]
が成り立つことである。

(1) より極大値は $f(0)=4>0$ であり，これは任意の $a>0$ に対して成り立つ。
ゆえに，極小値 $f(2a)=4-4a^{3}<0$ となればよい。
\[
  4-4a^{3}<0 \iff a^{3}-1>0 \iff (a-1)(a^{2}+a+1)>0
\]
ここで $a$ は実数であるから $a^{2}+a+1=\left(a+\dfrac{1}{2}\right)^{2}+\dfrac{3}{4}>0$ である。
よって $a-1>0$ すなわち $a>1$ を得る。これは $a>0$ を満たす。

\Ans{a>1}

\hs{確かめ}$a=1$ のときは $f(2)=0$ となって $x=2$ で $x$ 軸に接し，実数解は $x=-1,\ 2$ の2個にとどまる。
3個になるのは極小値が負に沈む $a>1$ のときである。

\Q{3}{10}
(1) の増減表より，$f(x)$ は区間 $0\leqq x\leqq2a$ で単調に減少し，
区間 $x\geqq2a$ で単調に増加する。
極小をとる $x=2a$（$>0$）が考える区間 $0\leqq x\leqq3$ の内部にあるかどうかで場合分けする。

\medskip
\noindent\textbf{［1］$0<2a<3$ すなわち $0<a<\dfrac{3}{2}$ のとき}

区間 $0\leqq x\leqq3$ において，$f(x)$ は $0\leqq x\leqq2a$ で減少し，$2a\leqq x\leqq3$ で増加する。
したがって，$f(x)$ は $x=2a$ で最小となり，最小値は
\[
  m(a)=f(2a)=4-4a^{3}
\]
である。

\medskip
\noindent\textbf{［2］$2a\geqq3$ すなわち $a\geqq\dfrac{3}{2}$ のとき}

区間 $0\leqq x\leqq3$ においてはつねに $0\leqq x\leqq3\leqq2a$ であるから，
$f'(x)\leqq0$ となり，$f(x)$ は $0\leqq x\leqq3$ 全体で単調に減少する。
したがって，$f(x)$ は右端 $x=3$ で最小となり，最小値は
\[
  m(a)=f(3)=3^{3}-3a\cdot3^{2}+4=31-27a
\]
である。

\medskip
以上［1］，［2］をまとめて，求める最小値 $m(a)$ は次の通りである。

\Ans{m(a)=\begin{cases} 4-4a^{3} & \left(0<a<\dfrac{3}{2}\right) \\[3mm] 31-27a & \left(a\geqq\dfrac{3}{2}\right) \end{cases}}

\hs{確かめ}境界 $a=\dfrac{3}{2}$ において，上の式は $4-4\cdot\left(\dfrac{3}{2}\right)^{3}=4-\dfrac{27}{2}=-\dfrac{19}{2}$，
下の式は $31-27\cdot\dfrac{3}{2}=31-\dfrac{81}{2}=-\dfrac{19}{2}$ となって一致する。
場合分けの関数を求めたあとは，必ず境界で値がつながることを確かめておきたい。

\begin{betsu}[グラフの位置関係による視覚的確認]
区間 $0\leqq x\leqq3$ の左端 $x=0$ は極大点であるから，$x>0$ に進むと $f(x)$ はまず減少する。
したがって，最小値をとる候補は極小点 $x=2a$（区間内にある場合）か右端 $x=3$ のいずれかに限られ，
左端 $x=0$ が最小になることはない。このことを答案の冒頭で述べておくと，端点 $f(0)=4$ との大小比較が不要である理由がより明確になる。
\end{betsu}

\end{document}
